\documentclass[11pt]{article}
\usepackage{amsmath,amssymb}
\usepackage[margin=2.3cm]{geometry}

\title{handcalcs al estilo Jupyter: \texttt{\%\%render} / \texttt{\%\%tex}}
\author{}
\date{\today}

\begin{document}
\maketitle

En Calc puedes usar \texttt{handcalcs} igual que en Jupyter: escribe una celda
que empiece (o contenga) la línea mágica \verb|%%render| o \verb|%%tex| y la
celda se \emph{renderiza automáticamente dentro del documento}. No se muestra
salida en la celda: el cálculo (fórmula, sustitución y resultado) aparece aquí,
justo donde está la celda.

\section{Tu ejemplo}
Exactamente el código que usabas en Jupyter:

%#python
import handcalcs.render
from math import sqrt, pi
%%tex
a = 2 / 3 * sqrt(pi)
%#end

\section{Cálculo de varias líneas}
Todo lo que va tras \verb|%%render| se renderiza línea a línea:

%#python
%%render
b = 20
h = 30
A = b * h
I = b * h**3 / 12
%#end

\section{Formato compacto (\texttt{params})}
La línea mágica admite los mismos argumentos que en Jupyter
(\verb|params|, \verb|long|, \verb|short|, precisión…):

%#python
%%render params 2
P = 1500
e = 35
M = P * e
%#end

\section{Con unidades (\texttt{pint})}
\texttt{pint} arrastra las unidades y \texttt{handcalcs} las muestra. El
registro \texttt{ureg} ya está disponible (no hace falta importar nada):

%#python
%%render
m = 1500 * ureg.kg
v = 25 * ureg.m / ureg.s
E_c = (m * v**2) / 2
%#end

\medskip
Cada bloque de arriba es una celda de Python que se ejecuta de verdad: las
variables siguen existiendo, así que también puedes usarlas en el texto con la
macro \verb|py|. Por ejemplo, el área de la sección era \py{A} y el momento
de inercia \py{f"{I:.0f}"}.

\end{document}
